
\section{Experiments}
%\subsection{Experimental setup}

In this section, we assess the efficacy of ProtoDiff in the context of three distinct few-shot learning scenarios: within-domain few-shot learning, cross-domain few-shot learning, and few-task few-shot learning.
%
For the within-domain few-shot learning experiments, we apply our method to three specific datasets: \textit{mini}Imagenet \cite{vinyals2016matching}, \textit{tiered}Imagenet \cite{ren2018meta}, and ImageNet-800 \cite{chen2021meta}.
%
Regarding cross-domain few-shot learning, we utilize \textit{mini}Imagenet \cite{vinyals2016matching} as the training domain, while testing is conducted on four distinct domains: CropDisease \cite{mohanty2016}, EuroSAT \cite{helber2019eurosat}, ISIC2018 \cite{tschandl2018ham10000}, and ChestX \cite{wang2017chestx}.
%
Furthermore, few-task few-shot learning~\cite{yao2021meta} challenges the common assumption of having abundant tasks available during meta-training. To explore this scenario, we perform experiments on four few-task meta-learning challenges: \textit{mini}Imagenet-S \cite{vinyals2016matching}, ISIC \cite{milton2019automated}, DermNet-S \cite{Dermnetdataset}, and Tabular Murris \cite{cao2020concept}.
%
For a comprehensive understanding of the datasets used in each setting, we provide detailed descriptions in the Appendix.


\textbf{Implementation details.} 
In our within-domain experiments, we utilize a Conv-4 and ResNet-12 backbone for \textit{mini}Imagenet and \textit{tiered}Imagenet. A ResNet-50 is used for ImageNet-800. We follow the approach described in \cite{chen2021meta} to achieve better performance and initially train a feature extractor on all the meta-training data without episodic training.  Standard data augmentation techniques are applied, including random resized crop and horizontal flip. 
%
For our cross-domain experiments, we use a ResNet-10 backbone to extract image features, which is a common choice for cross-domain few-shot classification \cite{ye20, guo2020broader}. The training configuration for this experiment is the same as the within-domain training. 
%
For few-task few-shot learning, we follow \cite{yao2021meta} using a network containing four convolutional blocks and a classifier layer. The average within-domain/ cross-domain, few-task few-shot classification accuracy ($\%$, top-1) along with $95\%$ confidence intervals are reported across all test query sets and tasks. The detailed implementation is provided in the appendix.


%\subsection{Results}


% \textbf{Visualization of the diffusion process.}
% The ProtoDiff method utilizes a task-guided diffusion model to generate prototypes that provide efficient class representations, as discussed in the previous section. To better understand the effectiveness of our proposed approach, we provide a visualization by Grad-Cam~\cite{selvaraju2017grad} in Figure~\ref{fig:diff_process} of the diffusion process, demonstrating how ProtoDiff gradually aggregates towards the desired class prototype during meta-training.  
% The vanilla prototype  does not exclusively focus on the face of the \textit{dog}, resulting in a lower probability.  ProtoDiff, with the diffusion process, randomly selects certain areas to add noise and perform diffusion, resulting in a prototype that gradually moves towards the face of the \textit{dog} with the highest probability at t=0. 
% The comparison between these different prototypes demonstrates the effectiveness of the ProtoDiff diffusion process in generating a more accurate and informative prototype for few-shot learning.



\textbf{Benefit of ProtoDiff.}
Table~\ref{tab:metadiff} compares our ProtoDiff with two baselines to demonstrate its effectiveness. The first Classifier-Baseline is a classification model trained on the entire label set using a classification loss and performing few-shot tasks with the cosine nearest-centroid method. The Meta-Baseline~\cite{chen2021meta} consists of two stages. In the first stage, a classifier is trained on all base classes, and the last fully connected layer is removed to obtain the feature encoder. The second stage is meta-learning, where the classifier is optimized using episodic training. The comparison between the Baseline and Meta-Baseline highlights the importance of episodic training for few-shot learning. ProtoDiff consistently outperforms Meta-Baseline~\cite{chen2021meta} by a large margin on all datasets and shots. The task-guided diffusion model employed in our ProtoDiff generates more informative prototypes, leading to improvements over deterministic prototypes.


\begin{table*}[t]
\caption{\textbf{Benefit of ProtoDiff} on  \textit{mini}Imagenet, \textit{tiered}Imagenet and ImageNet-800.   The results of the Classifier-Baseline and Meta-Baseline are provided by Chen~\etal~\cite{chen2021meta}. ProtoDiff consistently achieves better performance than two baselines on all datasets and settings.}
\centering
\scalebox{0.85}{
\begin{tabular}{lcccccc}
\toprule 
 &  \multicolumn{2}{c}{\textbf{\textit{mini}Imagenet}} & \multicolumn{2}{c}{\textbf{\textit{tiered}Imagenet}} & \multicolumn{2}{c}{\textbf{Imagenet-800}}  \\
\textbf{Method} & \textbf{1-shot} & \textbf{5-shot} & \textbf{1-shot} & \textbf{5-shot}  & \textbf{1-shot} & \textbf{5-shot}  \\
\midrule 

Classifier-Baseline~\cite{chen2021meta} & 58.91\textcolor{mygray}{\scriptsize{$\pm$0.23}} & 77.76\textcolor{mygray}{\scriptsize{$\pm$0.17}} & {68.07}\textcolor{mygray}{\scriptsize{$\pm$0.29}} & {83.74}\textcolor{mygray}{\scriptsize{$\pm$0.18}}  & {86.07}\textcolor{mygray}{\scriptsize{$\pm$0.21}} & {96.14}\textcolor{mygray}{\scriptsize{$\pm$0.08}} \\
Meta-Baseline~\cite{chen2021meta}  & 63.17\textcolor{mygray}{\scriptsize{$\pm$0.23}} & 79.26\textcolor{mygray}{\scriptsize{$\pm$0.17}} & {68.62}\textcolor{mygray}{\scriptsize{$\pm$0.27}} & {83.74}\textcolor{mygray}{\scriptsize{$\pm$0.18}} & {89.70}\textcolor{mygray}{\scriptsize{$\pm$0.19}} & {96.14}\textcolor{mygray}{\scriptsize{$\pm$0.08}} \\
\rowcolor{softblue}
{\textbf{ProtoDiff}}  & \textbf{66.63}\textcolor{mygray}{\scriptsize{$\pm$0.21}} & \textbf{83.48}\textcolor{mygray}{\scriptsize{$\pm$0.15}} & \textbf{72.95}\textcolor{mygray}{\scriptsize{$\pm$0.24}} &  \textbf{85.15}\textcolor{mygray}{\scriptsize{$\pm$0.18}} & \textbf{92.13}\textcolor{mygray}{\scriptsize{$\pm$0.20}} &  \textbf{98.21}\textcolor{mygray}{\scriptsize{$\pm$0.08}}
\\

\bottomrule
\end{tabular}}
\label{tab:metadiff}
\vspace{-5mm}
\end{table*}


\begin{figure*} [t!]
\centering
\includegraphics[width=0.99\linewidth]{graph/generative1.pdf}
\vspace{-2mm}
\caption{\textbf{Visualization of the different reconstructed overfitted prototypes using different generative models.} The vanilla prototype focuses only on the overall features of the bird and overlooks the finer details of the \textit{robin}. In contrast, the overfitted prototype can highlight specific features such as the tail and beak. While VAE and normalizing flow can generate prototypes of certain parts,  the diffusion method can generate prototypes that more closely resemble the overfitted prototype. With residual prototype learning, ProtoDiff achieves better.}
 \label{fig:generative}
 \vspace{-7mm}
\end{figure*}






\textbf{Effect of the residual prototype.}
The incorporation of residual prototypes in ProtoDiff presents a 
viable approach for reducing computational costs. This is attributable to the fact that the residual 
\begin{wrapfigure}{r}{0.45\linewidth}
\vspace{-6mm}
\centering
\begin{subfigure}[b]{0.45\textwidth}
  \centering
    \caption{Accuracy (\%) on various settings.}
    \scalebox{0.75}{
    \begin{tabular}[b]{lccc}
\toprule 
 &  \multicolumn{2}{c}{\textbf{\textit{mini}Imagenet}}  \\
\textbf{Method} & \textbf{1-shot} & \textbf{5-shot}  \\
\midrule 

Vanilla  & 63.17\textcolor{mygray}{\scriptsize{$\pm$0.23}} & 79.26\textcolor{mygray}{\scriptsize{$\pm$0.17}}  \\
ProtoDiff w/o residual & 64.75\textcolor{mygray}{\scriptsize{$\pm$0.22}} & 80.76\textcolor{mygray}{\scriptsize{$\pm$0.16}} \\
\rowcolor{softblue}
{\textbf{ProtoDiff w/ residual}}  & \textbf{66.63}\textcolor{mygray}{\scriptsize{$\pm$0.21}} & \textbf{83.48}\textcolor{mygray}{\scriptsize{$\pm$0.15}}
\\
\bottomrule
\end{tabular}}
  \label{table:residual}
\end{subfigure}
  \begin{subfigure}{0.4\textwidth}
  \centering
  \includegraphics[width=0.8\linewidth]{graph/loss_4.pdf}
  \vspace{-2mm}
  \caption{Optimization curves.}
  \label{fig:residual}
  \end{subfigure}
\caption{\textbf{Effect of residual prototype} on  \textit{mini}Imagenet. Utilizing ProtoDiff with a residual prototype not only accelerates the training process but also surpasses the performance of both the vanilla prototype and the non-residual prototype. }
  \vspace{-4mm}
\end{wrapfigure}
prototype exclusively encapsulates the disparity between the overfitted prototype and the original 
prototype, while the latter can be readily reused in each episode. Consequently, the calculation of the overfitted prototype is only necessitated once during the meta-training stage. Notably, the residual prototype often encompasses numerous zero values, thereby further expediting the training process. 
Table~\ref{table:residual} illustrates the superior accuracy achieved by ProtoDiff with residual prototype learning in comparison to both the vanilla model and ProtoDiff without the residual prototype. The training progress of ProtoDiff is visually represented in Figure~\ref{fig:residual}, demonstrating its accelerated training capabilities when compared to alternative methods. These results indicate the effectiveness of residual prototype learning in capturing task-specific information previously unaccounted for by the vanilla prototypes. Thus, the inclusion of residual prototype learning in ProtoDiff not only expedites the training process but also serves as a straightforward and effective approach to enhance the overall performance of ProtoDiff.


\textbf{Benefit of diffusion model.}
To confirm that the performance gain of our ProtoDiff model can be attributed to the diffusion model, we compared it with two widely used generative models: the variational autoencoder (VAE)~\cite{kingma2013auto} and normalizing flows~\cite{rezende2015variational}. VAE learns a low-dimensional representation of input data to generate new samples. Normalizing flows are more recent and learn a sequence of invertible transformations to map a simple prior distribution to the data distribution.
We obtained the task-specific prototype in the meta-test stage, conditioned on the support samples and $\epsilon \sim \mathcal{N}(0, I)$, by first acquiring an overfitted prototype $\mathbf{z}^*$ and then using it as the ground truth to train VAE and normalizing flows in the meta-training stage. The experimental results reported in Table~\ref{tab:ab_diff} show that the diffusion model outperforms both VAE and normalizing flows in terms of accuracy. Specifically, our diffusion model improves accuracy by 2.18\% compared to VAE and 1.52\% compared to normalizing flows.
Furthermore, we reconstructed the overfitted prototype using different generative models in Figure~\ref{fig:generative}. The vanilla prototype focuses solely on the bird's features and 
overlooks the finer details of the \textit{robin}. In contrast, the overfitted prototype can emphasize specific 
\begin{wraptable}{r}{7cm}
\vspace{-3mm}
\caption{\textbf{Benefit of diffusion model} over 
other generative models on  \textit{mini}Imagenet.}
\centering
\scalebox{0.85}{
\begin{tabular}{lcc}
\toprule 
 &  \multicolumn{2}{c}{\textbf{\textit{mini}Imagenet}}  \\
\textbf{Method} & \textbf{1-shot} & \textbf{5-shot}   \\
\midrule 
w/o generative model \cite{chen2021meta}  & 63.17\textcolor{mygray}{\scriptsize{$\pm$0.23}} & 79.26\textcolor{mygray}{\scriptsize{$\pm$0.17}} \\
w/ VAE & 64.45\textcolor{mygray}{\scriptsize{$\pm$0.22}} & 80.13\textcolor{mygray}{\scriptsize{$\pm$0.15}}  \\
w/ Normalizing flows & 65.11\textcolor{mygray}{\scriptsize{$\pm$0.22}} & 81.96\textcolor{mygray}{\scriptsize{$\pm$0.17}}  \\
\rowcolor{softblue}
{w/ \textbf{Diffusion}}  & \textbf{66.63}\textcolor{mygray}{\scriptsize{$\pm$0.21}} & \textbf{83.48}\textcolor{mygray}{\scriptsize{$\pm$0.15}} 
\\

\bottomrule
\end{tabular}}
\label{tab:ab_diff}
\end{wraptable}
features such as the tail and beak, resulting in better discrimination of the \textit{robin} from other bird species. While VAE and normalizing flows can generate prototypes of specific parts, our diffusion method can generate prototypes that more closely resemble the overfitted prototype, leading to improved performance.
Our findings suggest that diffusion models hold promise as a more practical approach for few-shot learning, as they can model complex distributions and produce informative prototypes.







\textbf{Visualization of the prototype adaptation process.}
We undertake a qualitative analysis of the adaptation process of our ProtoDiff framework 
during the meta-test time. We employ a 3-way 5-shot setup on the \textit{mini}ImageNet dataset for more comprehensive visualization of the adaptation process. To achieve this, we depict the vanilla prototypes, the diffused prototypes at various timesteps, and instances from the support and query sets. All examples are first normalized to a length of 1 and subsequently projected onto a two-dimensional space using t-SNE~\cite{van2008visualizing}. Figure~\ref{fig:tsne} displays the instance representations from both the query and support sets that are cross and circle, the vanilla prototypes that are denoted by a four-point star symbol, and the diffused prototypes that are denoted by stars.
\begin{wrapfigure}{r}{0.45\linewidth}
  \centering 
  \vspace{-3mm}
\includegraphics[width=0.9\linewidth] {graph/figure7-1.pdf}\\
    \caption{\textbf{t-SNE visualization of ProtoDiff adaptation process} on 3-way 5-shot task.}
    \label{fig:tsne}
    % \vspace{-8mm}
\end{wrapfigure}
Arrows depict the adaptation process of the diffused prototypes, and each color corresponds to each respective class.
Initially, the diffused prototypes are randomly allocated based on a Gaussian distribution, positioning them considerably from the optimal prototypes for each class. As the diffusion process continues, these diffused prototypes incrementally approach the optimal prototypes, resulting in more distinctive representations for different classes. In contrast, the vanilla prototypes lack distinguishability as they are computed using a single deterministic function based solely on the support set. This observation underscores the importance of the prototype adaptation process in attaining improved performance in few-shot classification tasks.








\textbf{Within-domain few-shot.}
We  evaluate our method on few-shot classification within domains,
in which the training domain is consistent with the test domain. The results are reported in Table~\ref{tab:main}. 
\begin{wraptable}{r}{8.5cm}
\vspace{-4mm}
\caption{\textbf{Within-domain few-shot} comparison on \textit{mini}Imagenet and \textit{tiered}Imagenet. ProtoDiff performs well on all datasets.}
\centering
\scalebox{0.65}{
\begin{tabular}{lcccc}
\toprule 
 &  \multicolumn{2}{c}{\textbf{\textit{mini}Imagenet}} & \multicolumn{2}{c}{\textbf{\textit{tiered}Imagenet}}  \\
\textbf{Method} & \textbf{1-shot} & \textbf{5-shot} & \textbf{1-shot} & \textbf{5-shot}  \\
\midrule 
TapNet \cite{yoon19}  & 61.65\textcolor{mygray}{\scriptsize{$\pm$0.15}} & 76.36\textcolor{mygray}{\scriptsize{$\pm$0.10}} & 63.08\textcolor{mygray}{\scriptsize{$\pm$0.15}} & 80.26\textcolor{mygray}{\scriptsize{$\pm$0.12}} \\
CTM \cite{li19}  & 62.05\textcolor{mygray}{\scriptsize{$\pm$0.55}} & 78.63\textcolor{mygray}{\scriptsize{$\pm$0.06}} & 64.78\textcolor{mygray}{\scriptsize{$\pm$0.11}} & 81.05\textcolor{mygray}{\scriptsize{$\pm$0.52}} \\
MetaOptNet \cite{lee19}  & 62.64\textcolor{mygray}{\scriptsize{$\pm$0.61}} & 78.63\textcolor{mygray}{\scriptsize{$\pm$0.46}} & 65.81\textcolor{mygray}{\scriptsize{$\pm$0.74}} & 81.75\textcolor{mygray}{\scriptsize{$\pm$0.53}} \\
Meta-Baseline \cite{chen2021meta}  & 63.17\textcolor{mygray}{\scriptsize{$\pm$0.23}} & 79.26\textcolor{mygray}{\scriptsize{$\pm$0.17}} & {68.62}\textcolor{mygray}{\scriptsize{$\pm$0.27}} & {83.74}\textcolor{mygray}{\scriptsize{$\pm$0.18}} \\
CAN \cite{hou19}  & 63.85\textcolor{mygray}{\scriptsize{$\pm$0.48}} & 79.44\textcolor{mygray}{\scriptsize{$\pm$0.34}} & {69.89}\textcolor{mygray}{\scriptsize{$\pm$0.51}} & {84.23}\textcolor{mygray}{\scriptsize{$\pm$0.37}} \\
Meta DeepBDC \cite{xie2022joint}  & 67.34\textcolor{mygray}{\scriptsize{$\pm$0.43}} & 84.46\textcolor{mygray}{\scriptsize{$\pm$0.28}} & {72.34}\textcolor{mygray}{\scriptsize{$\pm$0.49}} & {87.31}\textcolor{mygray}{\scriptsize{$\pm$0.32}} \\
SUN \cite{dong2022self}  & 67.80\textcolor{mygray}{\scriptsize{$\pm$0.45}} & 83.25\textcolor{mygray}{\scriptsize{$\pm$0.30}} & {72.99}\textcolor{mygray}{\scriptsize{$\pm$0.50}} & {86.74}\textcolor{mygray}{\scriptsize{$\pm$0.33}} \\
SetFeat~\cite{afrasiyabi2022matching} & 68.32\textcolor{mygray}{\scriptsize{$\pm$0.62}} & 82.71\textcolor{mygray}{\scriptsize{$\pm$0.46}} & {73.63}\textcolor{mygray}{\scriptsize{$\pm$0.88}} & {87.59}\textcolor{mygray}{\scriptsize{$\pm$0.57}} \\
\midrule
\rowcolor{softblue}
{\textbf{\textit{This paper}}} & \textbf{71.25}\textcolor{mygray}{\scriptsize{$\pm$0.45}} & \textbf{83.95}\textcolor{mygray}{\scriptsize{$\pm$0.45}} & \textbf{75.97}\textcolor{mygray}{\scriptsize{$\pm$0.75}} &  \textbf{88.75}\textcolor{mygray}{\scriptsize{$\pm$0.18}}
\\

\bottomrule
\vspace{-3mm}
\end{tabular}}
\label{tab:main}
\end{wraptable}
In this comparison, we apply SetFeat~\cite{afrasiyabi2022matching} with ProtoDiff to experiment since SetFeat is the current state-of-the-art model based on prototype-based meta-learning. Our ProtoDiff consistently achieves state-of-the-art performance on all datasets under various shots.
The better performance confirms that ProtoDiff can achieve more informative and higher quality prototypes to perform better for few-shot learning within domains.


 

\begin{table*}[t]
\centering
\caption{\textbf{Cross-domain few-shot} comparison on four benchmarks.  ProtoDiff is a consistent top-performer.}
\scalebox{0.8}{
\begin{tabular}{l c c c c c c c c}
    \toprule
&\multicolumn{2}{c}{\textbf{CropDiseases}}& \multicolumn{2}{c}{\textbf{EuroSAT}} & \multicolumn{2}{c}{\textbf{ISIC}}  & \multicolumn{2}{c}{\textbf{ChestX}}  \\
    %\cline{2-7}
    \textbf{Method} & \textbf{1-shot} &\textbf{5-shot} &\textbf{1-shot} &\textbf{5-shot}  &\textbf{1-shot} &\textbf{5-shot}  &\textbf{1-shot} &\textbf{5-shot} \\
    \midrule%\cline{2-7}

  
ProtoNet~\cite{snell2017prototypical}   &57.57\textcolor{mygray}{\scriptsize{$\pm$0.51}}  &{79.72}\textcolor{mygray}{\scriptsize{$\pm$0.67}} &{54.19}\textcolor{mygray}{\scriptsize{$\pm$0.57}} &{73.29}\textcolor{mygray}{\scriptsize{$\pm$0.71}}   &  29.62\textcolor{mygray}{\scriptsize{$\pm$0.32}}  & 39.57\textcolor{mygray}{\scriptsize{$\pm$0.57}} &{22.30}\textcolor{mygray}{\scriptsize{$\pm$0.25}}  &{24.05}\textcolor{mygray}{\scriptsize{$\pm$1.01}} \\

GNN \cite{satorras2018few}   &57.19\textcolor{mygray}{\scriptsize{$\pm$0.50}}  &{83.12}\textcolor{mygray}{\scriptsize{$\pm$0.40}} &{54.61}\textcolor{mygray}{\scriptsize{$\pm$0.50}} &{78.69}\textcolor{mygray}{\scriptsize{$\pm$0.40}}   &  30.14\textcolor{mygray}{\scriptsize{$\pm$0.30}}  & 42.54\textcolor{mygray}{\scriptsize{$\pm$0.40}}  &{21.94}\textcolor{mygray}{\scriptsize{$\pm$0.20}}  &{23.87}\textcolor{mygray}{\scriptsize{$\pm$0.20}} \\

AFA \cite{hu2022adversarial}   &67.61\textcolor{mygray}{\scriptsize{$\pm$0.50}}  &{88.06}\textcolor{mygray}{\scriptsize{$\pm$0.30}} &{63.12}\textcolor{mygray}{\scriptsize{$\pm$0.50}} &{85.58}\textcolor{mygray}{\scriptsize{$\pm$0.40}}   &  33.21\textcolor{mygray}{\scriptsize{$\pm$0.30}}  &\textbf {46.01}\textcolor{mygray}{\scriptsize{$\pm$0.40}}  &{22.92}\textcolor{mygray}{\scriptsize{$\pm$0.20}}  &{25.02}\textcolor{mygray}{\scriptsize{$\pm$0.20}} \\

ATA \cite{wang2021cross}   &67.47\textcolor{mygray}{\scriptsize{$\pm$0.50}}  &\textbf{90.59}\textcolor{mygray}{\scriptsize{$\pm$0.30}} &{61.35}\textcolor{mygray}{\scriptsize{$\pm$0.50}} &{83.75}\textcolor{mygray}{\scriptsize{$\pm$0.40}}   &  33.21\textcolor{mygray}{\scriptsize{$\pm$0.30}}  & 44.91\textcolor{mygray}{\scriptsize{$\pm$0.40}}  &{22.10}\textcolor{mygray}{\scriptsize{$\pm$0.20}}  &{24.32}\textcolor{mygray}{\scriptsize{$\pm$0.20}} \\

{HVM}~\cite{du2021hierarchical}  &{65.13}\textcolor{mygray}{\scriptsize{$\pm$0.45}}  &{87.65}\textcolor{mygray}{\scriptsize{$\pm$0.31}}  & {61.97}\textcolor{mygray}{\scriptsize{$\pm$0.34}}  &{74.88}\textcolor{mygray}{\scriptsize{$\pm$0.45}}  &{33.87}\textcolor{mygray}{\scriptsize{$\pm$0.35}}  &{42.05}\textcolor{mygray}{\scriptsize{$\pm$0.34}} & {\textbf{22.94}}\textcolor{mygray}{\scriptsize{$\pm$0.47}}  &{27.15}\textcolor{mygray}{\scriptsize{$\pm$0.45}} \\
\midrule
\rowcolor{softblue}
\textbf{\textit{This paper}}  &\textbf{68.93}\textcolor{mygray}{\scriptsize{$\pm$0.31}} & \textbf{90.15}\textcolor{mygray}{\scriptsize{$\pm$0.31}}  & \textbf{65.93}\textcolor{mygray}{\scriptsize{$\pm$0.34}} &\textbf{87.25}\textcolor{mygray}{\scriptsize{$\pm$0.35}}  &\textbf{34.97}\textcolor{mygray}{\scriptsize{$\pm$0.33}}   & \textbf{45.65}\textcolor{mygray}{\scriptsize{$\pm$0.31}}  & \textbf{23.01}\textcolor{mygray}{\scriptsize{$\pm$0.45}} &  \textbf{28.54}\textcolor{mygray}{\scriptsize{$\pm$0.41}}\\

   \bottomrule \\

\end{tabular}}
\label{tab:meta}
\vspace{-3mm}
\end{table*}
\textbf{Cross-domain few-shot.}
We also evaluate our method in the cross-domain few-shot classification, where the training domain is different from test domain.
Table~\ref{tab:meta} shows the evaluation of ProtoDiff on four datasets, with 5-way 1-shot and 5-way 5-shot settings. ProtoDiff achieves state-of-the-art performance on all four cross-domain few-shot learning benchmarks for each setting. On EuroSAT \cite{helber2019eurosat}, our model obtains high recognition accuracy under different shot configurations, outperforming the second best method, AFA~\cite{hu2022adversarial}, by a significant margin of $2.82\%$ for 5-way 1-shot. Even on the most challenging ChestX \cite{wang2017chestx}, which has a considerable domain gap with \textit{mini}ImageNet, our model achieves $28.54\%$ accuracy in the 5-way 5-shot setting, surpassing the second-best HVM~\cite{du2021hierarchical} by $1.39\%$. The consistent improvement across all benchmarks and settings confirms the effectiveness of ProtoDiff for cross-domain few-shot learning.

\textbf{Few-task few-shot.}
We evaluate ProtoDiff on the four different datasets under 5-way 1-shot and 5-way 5-shot in Table~\ref{tab:few-task}. In this comparison, we apply MLTI~\cite{yao2021meta} with ProtoDiff to experiment, since MLTI is the current state-of-the-art algorithm for few-task few-shot learning.
Our method achieves state-of-the-art performance on all four few-task meta-learning benchmarks under each setting.  On \textit{mini}Imagenet-S, our model achieves 44.15\% under 1-shot, surpassing the 
MLTI~\cite{yao2021meta}, by a margin of 2.79\%.   On ISIC~\cite{milton2019automated}, our method delivers 76.23\% for 5-shot, outperforming MLTI~\cite{yao2021meta} with 4.71\%. 
Even on the most challenging DermNet-S, which forms the largest dermatology dataset, our model delivers 51.45\% on the 5-way 1-shot setting. The consistent improvements on all benchmarks under various configurations confirm that our approach is also effective for few-task meta-learning.

\begin{table*}[t]
\centering
\caption{\textbf{Few-task few-shot} comparison on four benchmarks. ProtoDiff consistently ranks among the top-performing models.}
\scalebox{0.75}{
\begin{tabular}{lcccccccc}
\toprule
 & \multicolumn{2}{c}{\textbf{\textit{mini}Imagenet-S}} & \multicolumn{2}{c}{\textbf{ISIC}} & \multicolumn{2}{c}{\textbf{Dermnet-S}}  & \multicolumn{2}{c}{\textbf{Tabular Murris}}\\
& \textbf{1-shot} & \textbf{5-shot} & \textbf{1-shot} & \textbf{5-shot} & \textbf{1-shot} & \textbf{5-shot}  & \textbf{1-shot} & \textbf{5-shot}\\
\midrule
ProtoNet~\cite{snell2017prototypical}  & 36.26\textcolor{mygray}{\scriptsize{$\pm$0.70}} &  50.72\textcolor{mygray}{\scriptsize{$\pm$0.70}} & 58.56\textcolor{mygray}{\scriptsize{$\pm$1.01}} & 66.25\textcolor{mygray}{\scriptsize{$\pm$0.96}} & 44.21\textcolor{mygray}{\scriptsize{$\pm$0.75}} & 60.33\textcolor{mygray}{\scriptsize{$\pm$0.70}} & 80.03\textcolor{mygray}{\scriptsize{$\pm$0.90}} &  89.20 \textcolor{mygray}{\scriptsize{$\pm$0.56}} \\
Meta-Dropout~\cite{Lee2020Meta}  & 38.32\textcolor{mygray}{\scriptsize{$\pm$0.75}} &  52.53\textcolor{mygray}{\scriptsize{$\pm$0.70}} & 58.40\textcolor{mygray}{\scriptsize{$\pm$1.16}} & 67.32\textcolor{mygray}{\scriptsize{$\pm$0.95}} & 44.30\textcolor{mygray}{\scriptsize{$\pm$0.81}} & 60.86\textcolor{mygray}{\scriptsize{$\pm$0.75}} & 78.18\textcolor{mygray}{\scriptsize{$\pm$0.90}} &  89.25\textcolor{mygray}{\scriptsize{$\pm$0.56}} \\
MetaMix~\cite{yao2021improving} & 39.80\textcolor{mygray}{\scriptsize{$\pm$0.87}} & 53.35\textcolor{mygray}{\scriptsize{$\pm$0.88}} & 59.66\textcolor{mygray}{\scriptsize{$\pm$1.10}} & 68.97\textcolor{mygray}{\scriptsize{$\pm$0.80}} & 46.06\textcolor{mygray}{\scriptsize{$\pm$0.85}} & 62.97\textcolor{mygray}{\scriptsize{$\pm$0.70}} & 79.56\textcolor{mygray}{\scriptsize{$\pm$0.91}} & 88.88\textcolor{mygray}{\scriptsize{$\pm$0.60}} \\
Meta Interpolation~\cite{lee2022set} & {40.28}\textcolor{mygray}{\scriptsize{$\pm$0.70}} & {53.06}\textcolor{mygray}{\scriptsize{$\pm$0.72}} & - & - & - & -  & - & -\\
MLTI~\cite{yao2021meta}  & {41.36}\textcolor{mygray}{\scriptsize{$\pm$0.73}}  & {55.34}\textcolor{mygray}{\scriptsize{$\pm$0.72}} & {62.82}\textcolor{mygray}{\scriptsize{$\pm$1.09}} & {71.52}\textcolor{mygray}{\scriptsize{$\pm$0.89}} & {49.38}\textcolor{mygray}{\scriptsize{$\pm$0.85}} & {65.19}\textcolor{mygray}{\scriptsize{$\pm$0.73}} &81.89\textcolor{mygray}{\scriptsize{$\pm$0.88}} & 90.12\textcolor{mygray}{\scriptsize{$\pm$0.59}} \\
% MetaModulation~\cite{sun23metamodulation}  & {43.21}\textcolor{mygray}{\scriptsize{$\pm$0.73}}  & {57.26}\textcolor{mygray}{\scriptsize{$\pm$0.72}} & {65.61}\textcolor{mygray}{\scriptsize{$\pm$1.09}} & \textbf{76.40}\textcolor{mygray}{\scriptsize{$\pm$0.89}} & {50.45}\textcolor{mygray}{\scriptsize{$\pm$0.84}} & {67.05}\textcolor{mygray}{\scriptsize{$\pm$0.73}} &83.13\textcolor{mygray}{\scriptsize{$\pm$0.89}} & 91.23\textcolor{mygray}{\scriptsize{$\pm$0.57}} \\
\cmidrule{1-9}
\rowcolor{softblue}
\textbf{\textit{This paper}} & \textbf{44.75}\textcolor{mygray}{\scriptsize{$\pm$0.70}}  & \textbf{58.18}\textcolor{mygray}{\scriptsize{$\pm$0.72}}  & \textbf{66.13}\textcolor{mygray}{\scriptsize{$\pm$1.04}}   & {76.23}\textcolor{mygray}{\scriptsize{$\pm$0.81}}  &\textbf{51.53}\textcolor{mygray}{\scriptsize{$\pm$0.85}}  & \textbf{67.97}\textcolor{mygray}{\scriptsize{$\pm$0.74}}  &\textbf{84.03}\textcolor{mygray}{\scriptsize{$\pm$0.89}}    & \textbf{91.35}\textcolor{mygray}{\scriptsize{$\pm$0.57}}  \\
\bottomrule
\end{tabular}}
\vspace{-6mm}
\label{tab:few-task}
\end{table*}

\textbf{Limitations.} Naturally, our proposal also comes with limitations. Firstly, the diffusion model employed in our approach necessitates a substantial number of timesteps to sample the diffused prototype during the meta-test stage. Although we somewhat alleviate this issue by utilizing DDIM~\cite{song2020denoising} sampling, it still requires more computational resources than vanilla prototypes.
Secondly, obtaining the overfitted prototype involves fine-tuning, which inevitably leads to increased training time. While this fine-tuning step contributes to the effectiveness of our model, it comes with an additional computational cost.
As part of future work, we will investigate and address these limitations to further enhance the applicability and efficiency of our approach.